\subsection{射影线性簇}

设 W 是向量空间 V 的一个向量子空间； \(W - \{0\}\) 在由 V 导出的射影空间 \(\mathbf{P}(V)\) 中的典范像称为射影线性簇（若无混淆之虞，则简称为线性簇）；由于等价关系 \(\Delta(W)\) （定义在 \(W - \{0\}\) 上）是由关系 \(\Delta(V)\) 诱导的， \(W - \{0\}\) 在 \(\mathbf{P}(V)\) 中的像这一射影线性簇可以与由 W 导出的射影空间 \(\mathbf{P}(W)\) 等同起来，因而我们可以谈论这种簇的维数。在射影空间 \(\mathbf{P}(V)\) 中，V 的（去掉原点的）超平面的典范像是一个线性簇，称为射影超平面（或简称超平面）；若 \(\mathbf{P}(V)\) 是 n 维有限维的，则 \(\mathbf{P}(V)\) 中的超平面是维数为 n - 1 的线性簇。

关于向量空间的向量子空间的每一个命题，都可以转化为关于射影线性簇的命题。例如，如果一个射影空间 \(\mathbf{P}(\mathbf{V})\) 是有限维的，维数为 n，且 \((e_{i})_{0\leqslant i\leqslant n}\) 是 V 的一组基，那么每一个维数为 r 的线性簇 \(\mathbf{L}\subset\mathbf{P}(\mathbf{V})\) 都可以由 n-r 个齐次线性方程组成的方程组来定义

\[
\sum_ {i = 0} ^ {n} \xi_ {i} \alpha_ {i j} = 0 \quad (1 \leqslant j \leqslant n - r)\tag{4}
\]

在齐次坐标 \(\xi_{i}\) ( \(0 \leqslant i \leqslant n\) )（即 \(\mathbf{P}(\mathbf{V})\) 的点关于基 \((e_{i})\) 的齐次坐标）之间，其中 (4) 的左边是 V 上线性无关的形式。特别地，一个射影超平面由系数不全为零的单个齐次线性方程定义。反之， \(\mathbf{P}(\mathbf{V})\) 中满足关于 \(\xi_{i}\) 的任意齐次线性方程组的点构成一个线性簇 L；若所考虑的方程组由 \(k \leqslant n + 1\) 个方程组成，则 L 的维数 \(\geqslant n - k\) 。

 \(\mathbf{P}(\mathbf{V})\) 中线性簇的任何交都是一个线性簇；对 \(\mathbf{P}(\mathbf{V})\) 的每个子集 A，存在包含 A 的最小线性簇 L；它称为由 A 生成的线性簇，而 A 称为 L 的一个生成系。

如果 W 是 V 的由 \(\pi^{1}(A)\) ，则 L = P(W)。

若 \(\mathbf{L}\) 和 \(\mathbf{M}\) 是 \(\mathbf{P}(\mathbf{V})\) 中任意两个线性簇， \(\mathbf{N}\) 是由 \(\mathbf{L} \cup \mathbf{M}\) 生成的簇，则（§7，第 3 号，命题 4 的推论 3）

\[
\dim \mathbf {L} + \dim \mathbf {M} = \dim (\mathbf {L} \cap \mathbf {M}) + \dim \mathbf {N}.\tag{5}
\]

特别是，如果 \(\mathbf{P}(\mathbf{V})\) 是有限维的，且 \(\dim \mathbf{L} + \dim \mathbf{M} \geqslant \dim \mathbf{P}(\mathbf{V})\) ，由（5）可知 \(\mathbf{L} \cap \mathbf{M}\) 是非空的。

设 \((x_{i}), (y_{i})\) 是向量空间 V 中具有同一指标集的两族点，满足 \(y_{i} = \lambda_{i}x_{i}\) ，其中对所有 i 有 \(\lambda_{i}^{*} \neq 0\) 。若族 \((x_{i})\) 是自由的，则 \((y_{i})\) 也是自由的，反之亦然；这时称 P(V) 的点族 \(\pi(x_{i})\) 是射影自由的（或简称自由的）。这等价于说：对每个指标 K，点 \(\pi(x_{K})\) 不属于由诸 \(\pi(x_{i})\) （ \(i \neq K\) ）生成的线性簇。P(V) 中不是射影自由的点族称为射影相关的（或简称相关的）。

要使点的族 \((x_{i})\) （属于 \(\mathbf{V} - \{0\}\) ）满足：族 \((\pi (x_{i}))\) 射影自由且生成 \(\mathbf{P}(\mathbf{V})\) ，必要且充分条件是 \((x_{i})\) 是 \(\mathbf{V}\) 的一个基。若 \(\mathbf{P}(\mathbf{V})\) 的维数为 \(n\) ，则这样一个族中元素的个数是 \(n + 1\) 。注意，给定这样一个族 \((\pi (x_{i}))\) （在 \(\mathbf{P}(\mathbf{V})\) 中）并不确定 \(\mathbf{P}(\mathbf{V})\) 的一个给定点关于基 \((y_{i})\) （属于 \(\mathbf{V}\) ）的齐次坐标（即使允许相差一个左因子），其中 \(\pi (y_{i}) = \pi (x_{i})\) 对所有 \(\iota\) 成立（参见第 6 号）。

